% ==================================================
% CHAPTER 3: SYSTEM DESIGN
% ==================================================

\chapter{System Design}
\label{ch:design}

\section{Architecture and Decisions}

MemoryLane has three client-facing parts and one server
(Figure~\ref{fig:components}). Table~\ref{tab:responsibilities} gives each
part's responsibility and the requirements it serves.

\begin{table}[H]
\centering
\footnotesize
\begin{tabularx}{\textwidth}{>{\raggedright\arraybackslash}p{0.17\textwidth} L L >{\ttfamily}p{0.15\textwidth}}
\toprule
\textbf{Part} & \textbf{Responsibility} & \textbf{Technology} & \textrm{\textbf{Requirements}} \\
\midrule
Chrome extension & Capture pages and videos; automatic capture rules; sign-in & JavaScript, Manifest V3 & FR-02--FR-05 \\
Web app & Search, results, summaries, library, account & Svelte, Vite & FR-07--FR-10 \\
Caddy & HTTPS, serve the web app, forward \texttt{/api} & Caddy 2 & QR-05, QR-06 \\
FastAPI backend & Verify sign-in; save, search, summarise, report & Python, FastAPI, Sentence Transformers & FR-01--FR-12 \\
Apache Solr & Store chunks; BM25 and vector search & Solr 9.7 & FR-07, QR-01, QR-02 \\
SQLite & Sign-in links; problem reports; device keys & SQLite 3 & FR-01, FR-11, FR-15 \\
Local engine (extension) & Save, store (encrypted) and search pages on the device & Pyodide, transformers.js, IndexedDB & FR-13--FR-15, QR-12 \\
\bottomrule
\end{tabularx}
\caption{Architectural responsibilities}
\label{tab:responsibilities}
\end{table}

The design keeps the number of moving parts small enough for one server: the
backend is one process that loads the embedding model once, and Solr is the
only store for pages. Every important choice was recorded as an
architecture decision record (ADR) \cite{nygard2011adr}, with the
alternatives that were rejected and the cost of the choice
(Table~\ref{tab:adrs}).

\begin{table}[H]
\centering
\footnotesize
\begin{tabularx}{\textwidth}{>{\ttfamily}p{0.05\textwidth} >{\raggedright\arraybackslash}p{0.3\textwidth} >{\raggedright\arraybackslash}p{0.25\textwidth} L}
\toprule
\textrm{\textbf{\#}} & \textbf{Decision} & \textbf{Alternative rejected} & \textbf{Reason and trade-off} \\
\midrule
0001 & One Solr core for keyword and vector search & ChromaDB plus a separate BM25 index & One store to keep consistent; Solr is heavier to run \\
0002 & Fuse rankings with Reciprocal Rank Fusion & Min-max score averaging & No score normalisation or outlier effects; scores lose their scale \\
0003 & Drop weak matches with a relevance floor & Always show ten results & KNN always returns neighbours; a short list may look empty \\
0004 & Run everything on one EC2 instance & Managed services, several servers & Cost and simplicity; one point of failure \\
0005 & Link sign-ins by verified email & One library per sign-in & One library per person; needs a Cognito lookup \\
0006 & Fetch YouTube captions in the browser & Fetch on the server & YouTube blocks servers; depends on the user's tab \\
0007 & Store problem reports before emailing & Send email directly & Survives an email outage; one more database \\
0008 & Keep summaries on demand and optional & Summarise every search & Search never waits on Gemini; summaries cost a click \\
0009 & Remove the library map & Keep and improve it & Costly, fragile, rarely used \\
0010 & Remove ``Export my data'' & Keep it & Little use; deletion is the stronger right \\
0011 & Auto-capture keeps only pages that read like content & A list of blocked sites & Rules generalise; some content may be skipped \\
0012 & Pages on the device by default, server as fallback & Server only; local only & Server never holds them; device pages stay in one browser \\
0013 & Score vectors exactly & HNSW \texttt{\{!knn\}} & Both engines agree; 7\,ms against 5\,ms \\
0014 & Own BM25 and an int8 scan, no Orama & Orama & Half the latency, a tenth of the memory \\
0015 & The server's saving code in Pyodide & Readability in JavaScript & Identical chunks; 18\,MB runtime \\
0016 & Encrypt device pages; recovery-key backups & Plain storage; device-key backups & Server can't read backups; a lost key loses them \\
\bottomrule
\end{tabularx}
\caption{Architecture decision records}
\label{tab:adrs}
\end{table}

\section{Class Diagrams}

Figure~\ref{fig:class-backend} shows the core of the backend: the save
pipeline (\texttt{ingest}), the search (\texttt{retrieval}) and the only class
that talks to Solr, which adds the user filter to every request.
Figure~\ref{fig:class-client} shows the client side of deleting and signing
in, and Figure~\ref{fig:class-engine} the extension's local engine. The engine
depends on a \texttt{Storage} interface, not on IndexedDB: the extension uses
\texttt{IdbStorage}, the tests \texttt{MemoryStorage}. \texttt{LocalIndex}
answers the same questions the server's \texttt{SolrClient} does, so the
ported search runs unchanged over either. Python and JavaScript modules are
drawn as classes with the stereotype $\ll$module$\gg$; dashed arrows are
dependencies, filled diamonds composition, open triangles realisation.

\begin{figure}[H]
\centering
\begin{tikzpicture}[>={Stealth[length=2mm]}]
  % top row
  \node[umlclass=30mm, anchor=north] (error) at (-5.55,0) {%
    \centering$\ll$exception$\gg$\\\textbf{IngestError}
    \nodepart{second}+ status\_code: int\\+ detail: str
    \nodepart{third}+ \_\_init\_\_(status\_code,\\ \quad detail)};
  \node[umlclass=58mm, anchor=north] (ingest) at (-0.2,0) {%
    \centering$\ll$module$\gg$\\\textbf{ingest}
    \nodepart{second}-- \_user\_write\_lock(user\_id): RLock
    \nodepart{third}+ ingest\_page(user\_id, url, title, text,\\ \quad html, canonical\_url, is\_manual,\\ \quad keep\_existing, youtube): IngestResult\\
      + resolve\_canonical\_url(url, canonical): str\\+ page\_vector(embeddings): list};
  \node[umlclass=48mm, anchor=north] (retrieval) at (5.65,0) {%
    \centering$\ll$module$\gg$\\\textbf{retrieval}
    \nodepart{second}RRF\_K, MIN\_VECTOR\_SIMILARITY
    \nodepart{third}+ hybrid\_search(query, user\_id, top\_k)\\+ rrf\_fuse\_lists(lists, k): dict\\+ group\_by\_webpage(chunks, scores)\\+ apply\_relevance\_floor(pages)};
  % bottom row
  \node[umlclass=30mm, anchor=north] (result) at (-5.55,-4.75) {%
    \centering$\ll$dataclass$\gg$\\\textbf{IngestResult}
    \nodepart{second}+ status: str\\+ message: str\\+ page\_id: str\\+ url, title, domain: str\\+ timestamp: str\\+ chunks: int
    \nodepart{third}\textit{(none)}};
  \node[umlclass=58mm, anchor=north] (solr) at (-0.2,-4.75) {%
    \textbf{SolrClient}
    \nodepart{second}+ base\_url: str\\+ client: httpx.Client
    \nodepart{third}+ upsert\_page\_chunks(page\_id, docs, user\_id)\\+ delete\_page(page\_id, user\_id)\\+ bm25\_search(query, user\_id, rows): list\\+ knn\_search(vector, user\_id, k): list\\+ get\_page\_head(page\_id, user\_id): dict\\+ list\_pages(user\_id, limit, offset)\\+ existing\_page\_ids(user\_id, ids): set\\-- \_update(body, visible\_now)};
  \node[umlclass=48mm, anchor=north] (intent) at (5.65,-4.75) {%
    \centering$\ll$dataclass$\gg$\\\textbf{QueryIntent}
    \nodepart{second}+ text: str\\+ domains: set[str]\\+ saved\_from, saved\_until: datetime\\+ saved\_near: datetime\\+ url: str
    \nodepart{third}+ active(): bool};

  \draw[->, densely dashed] (ingest.west |- error) -- node[font=\tiny, above] {$\ll$raise$\gg$} (error.east);
  \draw[->, densely dashed] (ingest.south west) -- node[font=\tiny, above, sloped] {$\ll$create$\gg$} (result.north east);
  \draw[->, densely dashed] (ingest.south) -- node[font=\tiny, right] {$\ll$use$\gg$} (solr.north);
  \draw[->, densely dashed] (retrieval.south west) -- node[font=\tiny, above, sloped] {$\ll$use$\gg$} (solr.north east);
  \draw[->, densely dashed] (retrieval.south) -- node[font=\tiny, right] {$\ll$use$\gg$} (intent.north);
\end{tikzpicture}

\caption{Class diagram: backend save and search}
\label{fig:class-backend}
\end{figure}

\begin{figure}[H]
\centering
\begin{tikzpicture}[>={Stealth[length=2mm]}]
  \node[umlclass=50mm] (queue) at (0,0) {%
    \textbf{DeleteQueue}
    \nodepart{second}+ hidden: Store\textless Set\textless id\textgreater\textgreater\\+ toasts: Store\textless Toast[]\textgreater\\-- pending: Map\textless id, timer\textgreater\\-- committedAt: Map\textless id, time\textgreater
    \nodepart{third}+ schedule(pageId)\\+ undo(pageId)\\+ flush(\{keepalive\}): int\\+ restore(): Promise\textless int\textgreater\\+ reconcile(ids, fetchedAt)\\+ notInTotal(id, fetchedAt): bool};
  \node[umlclass=46mm, anchor=north west] (api) at (3.4,1.6) {%
    \centering$\ll$module$\gg$\\\textbf{api}
    \nodepart{second}\textit{API\_BASE, TIMEOUT\_MS}
    \nodepart{third}+ search(query, signal)\\+ listMemories(limit, offset)\\+ fetchSummary(\{query, id\})\\+ deletePage(id, \{keepalive\})\\+ deleteAllData()};
  \node[umlclass=30mm, anchor=north west] (apierr) at (9.1,1.6) {%
    \textbf{ApiError}
    \nodepart{second}+ message: str\\+ status: int
    \nodepart{third}\textit{extends Error}};
  \node[umlclass=46mm, anchor=north west] (auth) at (3.4,-2.2) {%
    \centering$\ll$module$\gg$\\\textbf{auth}
    \nodepart{second}+ session: Store\textless\{email\}\textgreater\\-- access: \{token, expiresAt\}
    \nodepart{third}+ signIn(\{provider, signUp\})\\+ handleCallback()\\+ getAccessToken(\{force\}): token\\+ signOut()\\+ refreshRejected(status): bool};
  \node[umlclass=30mm, anchor=north west] (sesserr) at (9.1,-2.2) {%
    \textbf{SessionUnavail-\\ableError}
    \nodepart{second}\textit{(none)}
    \nodepart{third}\textit{extends Error}};

  \draw[->, densely dashed] (queue.east |- api.west) -- node[font=\tiny, above] {$\ll$use$\gg$} (api.west);
  \draw[->, densely dashed] (api.south) -- node[font=\tiny, right] {$\ll$use$\gg$} (auth.north);
  \draw[->, densely dashed] (api.east |- apierr.west) -- node[font=\tiny, above] {$\ll$throw$\gg$} (apierr.west);
  \draw[->, densely dashed] (auth.east |- sesserr.west) -- node[font=\tiny, above] {$\ll$throw$\gg$} (sesserr.west);
\end{tikzpicture}

\caption{Class diagram: web app deletes and sessions}
\label{fig:class-client}
\end{figure}

\begin{figure}[H]
\centering
\begin{tikzpicture}[>={Stealth[length=2mm]}]
  % three columns, 44 mm wide, 6 mm apart
  \def\ca{0} \def\cb{5.0} \def\cc{10.0}
  % ---- row 1: the engine and what it uses --------------------------------------
  \node[umlclass=40mm, anchor=north west] (emb) at (\ca,0) {%
    \centering$\ll$module$\gg$\\\textbf{embedder}
    \nodepart{second}bge-small; WebGPU or WASM
    \nodepart{third}+ embedQuery(q)\\+ embedPassages(texts)};
  \node[umlclass=44mm, anchor=north west] (engine) at (\cb,0) {%
    \textbf{LocalEngine}
    \nodepart{second}+ libraryId: str\\-- queue: one write at a time
    \nodepart{third}+ open(): stats\\+ save(body): \{status, page\}\\+ search(query, n): results\\+ status(url, canonical)\\+ list(limit, offset)\\+ exportPage(id) / importPage(p)\\+ deletePage(id)};
  \node[umlclass=40mm, anchor=north west] (ing) at (\cc,0) {%
    \centering$\ll$module$\gg$\\\textbf{ingest}
    \nodepart{second}Pyodide + backend/ingest.py
    \nodepart{third}+ preparePage(request,\\\quad lookupHead)};
  % ---- row 2: what the engine is made of ---------------------------------------
  \node[umlclass=40mm, anchor=north west] (store) at (\ca,-4.9) {%
    \textbf{PageStore}
    \nodepart{second}-- key: CryptoKey (device key)\\-- sizes: Map\textless id, bytes\textgreater
    \nodepart{third}+ loadAll(): pages\\+ putPage(page, chunks)\\+ readPage(id) / deletePage(id)};
  \node[umlclass=44mm, anchor=north west] (index) at (\cb,-4.9) {%
    \textbf{LocalIndex}
    \nodepart{second}-- chunks, pages: Map\\-- vectors: int8 + scale
    \nodepart{third}+ bm25Search(q, \{rows, filters\})\\+ knnSearch(v, \{k, pageIds\})\\+ recentPages(f) / listDomains()\\+ getPageHead(id)};
  \node[umlclass=40mm, anchor=north west] (bm25) at (\cc,-4.9) {%
    \textbf{Bm25Index}
    \nodepart{second}-- title, text: postings, norms
    \nodepart{third}+ add(doc, title, text)\\+ remove(doc)};
  % ---- row 3: storage behind an interface --------------------------------------
  \node[umlclass=40mm, anchor=north west] (storage) at (\ca,-8.6) {%
    \centering$\ll$interface$\gg$\\\textbf{Storage}
    \nodepart{second}\textit{(none)}
    \nodepart{third}+ getAll() / get(id)\\+ put(record) / delete(id)};
  \node[umlclass=44mm, anchor=north west] (idb) at (\cb,-8.6) {%
    \textbf{IdbStorage}
    \nodepart{second}-- db: one IndexedDB per library
    \nodepart{third}+ destroy()};
  \node[umlclass=40mm, anchor=north west] (mem) at (\cc,-8.6) {%
    \textbf{MemoryStorage}
    \nodepart{second}-- records: Map
    \nodepart{third}\textit{for tests}};

  % ---- relationships -----------------------------------------------------------
  \draw[->, densely dashed] (engine.west |- emb.text split) -- (emb.east |- emb.text split);
  \draw[->, densely dashed] (engine.east |- ing.text split) -- (ing.west |- ing.text split);
  \draw[-{Diamond[fill=black, length=2.6mm]}] (store.north) -- ++(0,0.35) -| ([xshift=-6mm]engine.south);
  \draw[-{Diamond[fill=black, length=2.6mm]}] (index.north) -- (index.north |- engine.south);
  \draw[-{Diamond[fill=black, length=2.6mm]}] (bm25.west |- index.text split) -- (index.east |- index.text split);
  \draw[->, densely dashed] (store.south) -- node[font=\tiny, right] {$\ll$use$\gg$} (storage.north);
  \draw[-{Triangle[open, length=2.6mm]}, densely dashed] (idb.west |- storage.text split) -- (storage.east |- storage.text split);
  \draw[-{Triangle[open, length=2.6mm]}, densely dashed] (mem.south) -- ++(0,-0.35) -| ([xshift=6mm]storage.south);
\end{tikzpicture}

\caption{Class diagram: the local engine in the extension}
\label{fig:class-engine}
\end{figure}

\section{Sequence Diagrams}

Figure~\ref{fig:seq-save} shows a manual save to the server, including the
case where the page has not changed; Figure~\ref{fig:seq-save-device} the same
save kept on the device (the default), where the server's saving code runs in
Pyodide; Figure~\ref{fig:seq-search} a search, with the device and the server
searched in parallel and their results fused; and Figure~\ref{fig:seq-signin} a
sign-in with PKCE followed by the first API request, which finds the library
for the sign-in.

\begin{figure}[H]
\centering
\begin{tikzpicture}[x=2.55cm, y=0.5cm]
  \foreach \i/\t in {0/{Reader}, 1/{:Popup}, 2/{:Service\\worker}, 3/{:Page tab}, 4/{:FastAPI}, 5/{:Solr}} {
    \node[lifehead] (h\i) at (\i,0) {\t};
    \draw[life] (h\i.south) -- (\i,-18.6);
  }
  \node[act, fit={(2,-2.5) (2,-17.1)}] {};
  \node[act, fit={(4,-5.3) (4,-15.7)}] {};
  \newcommand{\m}[4]{\draw[msg] (#1,#3) -- node[mlabel, fill=white] {#4} (#2,#3);}
  \newcommand{\rt}[4]{\draw[ret] (#1,#3) -- node[mlabel, fill=white] {#4} (#2,#3);}
  \newcommand{\self}[3]{\draw[msg] (#1,#2) -- ++(0.12,0) -- ++(0,-0.75) node[font=\scriptsize, right, pos=0.5, align=left, fill=white, inner sep=1pt] {#3} -- ++(-0.12,0);}

  \m{0}{1}{-1.6}{Remember this page}
  \m{1}{2}{-2.5}{save(tab)}
  \m{2}{3}{-3.4}{inject content.js}
  \rt{3}{2}{-4.3}{HTML, text, canonical URL}
  \m{2}{4}{-5.3}{POST /api/memory (Bearer)}
  \self{4}{-5.9}{verify token\\$\rightarrow$ library id}
  \self{4}{-7.4}{normalise URL,\\page\_id; clean}
  \m{4}{5}{-9.1}{get\_page\_head}
  \rt{5}{4}{-10.0}{head or none}
  % alt frame
  \draw (2.55,-10.6) rectangle (5.4,-16.3);
  \node[draw, fill=white, font=\tiny, anchor=north west, inner sep=1.5pt] at (2.55,-10.6) {alt};
  \node[font=\tiny, anchor=west] at (2.75,-11.05) {[same text and model, or automatic save of a saved page]};
  \rt{4}{2}{-12.0}{unchanged}
  \draw[densely dashed] (2.55,-12.5) -- (5.4,-12.5);
  \node[font=\tiny, anchor=west] at (2.75,-12.95) {[new or changed]};
  \self{4}{-13.4}{chunk, embed}
  \m{4}{5}{-14.9}{upsert chunks}
  \rt{4}{2}{-15.7}{saved, page}
  \rt{2}{1}{-17.1}{result}
  \rt{1}{0}{-18.0}{``Saved just now'' + Undo}
\end{tikzpicture}

\caption{Sequence diagram: saving a page on the server}
\label{fig:seq-save}
\end{figure}

\begin{figure}[H]
\centering
\begin{tikzpicture}[x=2.1cm, y=0.5cm]
  \foreach \i/\t in {0/{Reader}, 1/{:Service\\worker}, 2/{:Local\\engine}, 3/{:Pyodide\\(ingest\_page)}, 4/{:Embedding\\model}, 5/{:IndexedDB}, 6/{:FastAPI}} {
    \node[lifehead] (h\i) at (\i,0) {\t};
    \draw[life] (h\i.south) -- (\i,-21.0);
  }
  \node[act, fit={(1,-1.6) (1,-20.2)}] {};
  \node[act, fit={(2,-2.5) (2,-15.4)}] {};
  \node[act, fit={(3,-3.4) (3,-9.2)}] {};
  \newcommand{\m}[4]{\draw[msg] (#1,#3) -- node[mlabel, fill=white] {#4} (#2,#3);}
  \newcommand{\rt}[4]{\draw[ret] (#1,#3) -- node[mlabel, fill=white] {#4} (#2,#3);}
  \newcommand{\self}[3]{\draw[msg] (#1,#2) -- ++(0.1,0) -- ++(0,-0.75) node[font=\scriptsize, right, pos=0.5, align=left, fill=white, inner sep=1pt] {#3} -- ++(-0.1,0);}

  \m{0}{1}{-1.6}{Remember this page}
  \m{1}{2}{-2.5}{SAVE (POST /memory body)}
  \m{2}{3}{-3.4}{ingest\_page(body)}
  \self{3}{-3.9}{canonical URL $\rightarrow$\\page\_id; clean (trafilatura)}
  \m{3}{2}{-5.6}{get\_page\_head(page\_id)}
  \rt{2}{3}{-6.5}{head or none}
  \self{3}{-7.0}{chunk; same refusals\\as the server}
  \rt{3}{2}{-9.2}{chunks, texts to embed}
  \m{2}{4}{-10.1}{embed ``title + chunk''}
  \rt{4}{2}{-11.0}{384-d vectors}
  \self{2}{-11.5}{keyword statistics;\\int8 vectors}
  \m{2}{5}{-13.6}{put encrypted record}
  \rt{2}{1}{-15.4}{saved, page}
  % opt frame: past the size limit
  \draw (0.6,-16.0) rectangle (6.4,-19.4);
  \node[draw, fill=white, font=\tiny, anchor=north west, inner sep=1.5pt] at (0.6,-16.0) {opt};
  \node[font=\tiny, anchor=west] at (0.9,-16.45) {[past the size limit, and the reader agreed]};
  \m{1}{6}{-17.4}{POST /api/memory/transfer (oldest pages, as chunks)}
  \rt{6}{1}{-18.6}{saved on the server; then removed here}
  \rt{1}{0}{-20.2}{``Saved just now'' + Undo}
\end{tikzpicture}

\caption{Sequence diagram: saving a page on the device}
\label{fig:seq-save-device}
\end{figure}

\begin{figure}[H]
\centering
\begin{tikzpicture}[x=2.55cm, y=0.5cm]
  \foreach \i/\t in {0/{Reader}, 1/{:Web app}, 2/{:Extension\\(local engine)}, 3/{:FastAPI}, 4/{:Embedding\\model}, 5/{:Solr}} {
    \node[lifehead] (h\i) at (\i,0) {\t};
    \draw[life] (h\i.south) -- (\i,-20.4);
  }
  \node[act, fit={(1,-1.6) (1,-19.8)}] {};
  \node[act, fit={(2,-3.4) (2,-5.0)}] {};
  \node[act, fit={(3,-6.6) (3,-17.0)}] {};
  \newcommand{\m}[4]{\draw[msg] (#1,#3) -- node[mlabel, fill=white] {#4} (#2,#3);}
  \newcommand{\rt}[4]{\draw[ret] (#1,#3) -- node[mlabel, fill=white] {#4} (#2,#3);}
  \newcommand{\self}[3]{\draw[msg] (#1,#2) -- ++(0.1,0) -- ++(0,-0.75) node[font=\scriptsize, right, pos=0.5, align=left, fill=white, inner sep=1pt] {#3} -- ++(-0.1,0);}

  \m{0}{1}{-1.6}{type or speak a query}
  % outer par: pages on this device, pages on the server
  \draw (0.55,-2.3) rectangle (5.4,-17.6);
  \node[draw, fill=white, font=\tiny, anchor=north west, inner sep=1.5pt] at (0.55,-2.3) {par};
  \node[font=\tiny, anchor=east] at (5.35,-2.75) {[pages on this device]};
  \m{1}{2}{-3.4}{SEARCH (this library only)}
  \self{2}{-3.9}{same pipeline,\\over IndexedDB}
  \rt{2}{1}{-5.0}{ranked pages}
  \draw[densely dashed] (0.55,-5.6) -- (5.4,-5.6);
  \node[font=\tiny, anchor=east] at (5.35,-6.0) {[pages on the server]};
  \m{1}{3}{-6.6}{POST /api/search \{query\}}
  \self{3}{-7.1}{verify token\\$\rightarrow$ library id}
  % inner par: keyword and vector search
  \draw (2.6,-8.5) rectangle (5.3,-13.4);
  \node[draw, fill=white, font=\tiny, anchor=north west, inner sep=1.5pt] at (2.6,-8.5) {par};
  \m{3}{5}{-9.2}{BM25, user filter, top 50}
  \draw[densely dashed] (2.6,-9.8) -- (5.3,-9.8);
  \m{3}{4}{-10.5}{embed query (cached)}
  \rt{4}{3}{-11.3}{384-d vector}
  \m{3}{5}{-12.2}{exact vectors, top 50}
  \rt{5}{3}{-13.0}{ranked chunk lists}
  \self{3}{-13.9}{RRF fusion;\\group into pages}
  \self{3}{-15.3}{relevance floor;\\snippets}
  \rt{3}{1}{-17.0}{up to 10 pages}
  \self{1}{-18.1}{fuse both lists\\(RRF, k = 60)}
  \rt{1}{0}{-19.8}{results with snippets}
\end{tikzpicture}

\caption{Sequence diagram: a search over the device and the server}
\label{fig:seq-search}
\end{figure}

\begin{figure}[H]
\centering
\begin{tikzpicture}[x=3.05cm, y=0.5cm]
  \foreach \i/\t in {0/{User}, 1/{:Web app}, 2/{:Amazon\\Cognito}, 3/{:FastAPI}, 4/{:accounts.db}} {
    \node[lifehead] (h\i) at (\i,0) {\t};
    \draw[life] (h\i.south) -- (\i,-19.0);
  }
  \newcommand{\m}[4]{\draw[msg] (#1,#3) -- node[mlabel, fill=white] {#4} (#2,#3);}
  \newcommand{\rt}[4]{\draw[ret] (#1,#3) -- node[mlabel, fill=white] {#4} (#2,#3);}
  \newcommand{\self}[3]{\draw[msg] (#1,#2) -- ++(0.1,0) -- ++(0,-0.75) node[font=\scriptsize, right, pos=0.5, align=left, fill=white, inner sep=1pt] {#3} -- ++(-0.1,0);}

  \m{0}{1}{-1.6}{Sign in}
  \self{1}{-2.2}{make verifier and\\SHA-256 challenge}
  \m{1}{2}{-3.9}{/oauth2/authorize (challenge)}
  \m{2}{0}{-4.8}{hosted sign-in (Google or email)}
  \m{0}{2}{-5.7}{credentials}
  \rt{2}{1}{-6.6}{redirect with one-time code}
  \m{1}{2}{-7.5}{/oauth2/token (code + verifier)}
  \rt{2}{1}{-8.4}{access + refresh tokens}
  \m{1}{3}{-9.5}{API request, Bearer access token}
  \self{3}{-10.1}{verify RS256, exp,\\iss, token\_use, client\_id}
  \m{3}{4}{-11.8}{library for this sub?}
  % alt frame
  \draw (2.55,-12.4) rectangle (4.35,-17.6);
  \node[draw, fill=white, font=\tiny, anchor=north west, inner sep=1.5pt] at (2.55,-12.4) {alt};
  \node[font=\tiny, anchor=west] at (2.6,-12.95) {[known sign-in]};
  \rt{4}{3}{-13.5}{library id}
  \draw[densely dashed] (2.55,-14.1) -- (4.35,-14.1);
  \node[font=\tiny, anchor=east] at (4.3,-14.55) {[first time]};
  \m{3}{2}{-15.2}{/oauth2/userInfo (user's token)}
  \rt{2}{3}{-16.0}{email, email\_verified}
  \m{3}{4}{-16.9}{link sub to library of verified email}
  \rt{3}{1}{-18.4}{response (only that library's data)}
\end{tikzpicture}

\caption{Sequence diagram: sign-in (authorization code with PKCE) and the first API request}
\label{fig:seq-signin}
\end{figure}

\section{Collaboration Diagrams}

Collaboration (communication) diagrams show the same kind of interaction as
links between objects, with numbered messages. Figure~\ref{fig:collab-search}
shows the objects inside the backend during a search;
Figure~\ref{fig:collab-delete} shows a delete with undo, where message 2a
(undo) and 2b (commit) are alternatives.

\begin{figure}[H]
\centering
\begin{tikzpicture}[>={Stealth[length=1.6mm]}]
  \node[obj] (web)  at (0, 0)     {\underline{:Web app}};
  \node[obj] (main) at (0,-2.0)   {\underline{:main /search}};
  \node[obj] (ret)  at (0,-4.0)   {\underline{:retrieval}};
  \node[obj] (int)  at (-4.6,-4.0) {\underline{:query\_intent}};
  \node[obj] (solr) at (4.8,-4.0) {\underline{:SolrClient}};
  \node[obj] (emb)  at (0,-6.0)   {\underline{:embeddings}};
  \node[obj] (snip) at (4.8,-2.0) {\underline{:snippets}};

  \foreach \a/\b in {web/main, main/ret, ret/int, ret/solr, ret/emb, main/snip}
    \draw (\a) -- (\b);
  \draw (ret.south west) .. controls ++(-0.9,-0.9) and ++(-0.9,0.4) .. (ret.west);

  \node[cmsg, right=2mm] at ($(web)!0.5!(main)$) {1: search(query) $\downarrow$};
  \node[cmsg, left=2mm]  at ($(web)!0.5!(main)$) {3: up to 10 pages $\uparrow$};
  \node[cmsg, right=2mm] at ($(main)!0.5!(ret)$) {1.1: hybrid\_search(q, uid) $\downarrow$};
  \node[cmsg, above=1.5mm] at ($(main)!0.5!(snip)$) {1.2: snippets(pages, q) $\rightarrow$};
  \node[cmsg, above=1.5mm] at ($(ret)!0.5!(int)$) {$\leftarrow$ 1.1.1: parse\_query(q, sites)};
  \node[cmsg, above=1.5mm] at ($(ret)!0.5!(solr)$) {1.1.2: bm25\_search(q, uid) $\rightarrow$};
  \node[cmsg, below=1.5mm] at ($(ret)!0.5!(solr)$) {1.1.4: knn\_search(v, uid) $\rightarrow$};
  \node[cmsg, right=2mm] at ($(ret)!0.5!(emb)$) {1.1.3: embed\_query(q) $\downarrow$};
  \node[cmsg, anchor=east] at ($(ret.west)+(-0.95,-0.55)$) {1.1.5: fuse, group,\\relevance floor};
\end{tikzpicture}

\caption{Collaboration diagram: a search inside the backend}
\label{fig:collab-search}
\end{figure}

\begin{figure}[H]
\centering
\begin{tikzpicture}[>={Stealth[length=1.6mm]}]
  \node[obj] (row)   at (-5.2, 0)   {\underline{:PageRow}};
  \node[obj] (queue) at (-0.9, 0)   {\underline{:DeleteQueue}};
  \node[obj] (toast) at (-0.9, 2.3) {\underline{:Toasts}};
  \node[obj] (api)   at (3.2, 0)    {\underline{:api}};
  \node[obj] (back)  at (3.2,-2.1)  {\underline{:FastAPI}};
  \node[obj] (solr)  at (3.2,-4.2)  {\underline{:SolrClient}};

  \foreach \a/\b in {row/queue, queue/toast, queue/api, api/back, back/solr}
    \draw (\a) -- (\b);
  \draw (queue.south west) .. controls ++(-0.5,-1.0) and ++(0.5,-1.0) .. (queue.south east);

  \node[cmsg, above=1.5mm] at ($(row)!0.5!(queue)$) {1: schedule(id) $\rightarrow$};
  \node[cmsg, anchor=north] at ($(queue.south)+(0,-1.0)$) {1.1: hide(id)\\2b: commit(id) [after 6\,s]};
  \node[cmsg, left=2mm]  at ($(queue)!0.5!(toast)$) {1.2: show ``Page deleted $\cdot$ Undo'' $\uparrow$};
  \node[cmsg, right=2mm] at ($(queue)!0.5!(toast)$) {$\downarrow$ 2a: undo(id) [within 6\,s]};
  \node[cmsg, above=1.5mm] at ($(queue)!0.5!(api)$) {2b.1: deletePage(id) $\rightarrow$};
  \node[cmsg, right=2mm] at ($(api)!0.5!(back)$) {2b.1.1: DELETE /api/memory/\{id\} $\downarrow$};
  \node[cmsg, right=2mm] at ($(back)!0.5!(solr)$) {2b.1.1.1: delete\_page(id, uid),\\soft commit $\downarrow$};
\end{tikzpicture}

\caption{Collaboration diagram: deleting a page with undo}
\label{fig:collab-delete}
\end{figure}

\section{State Chart Diagrams}

Figure~\ref{fig:state-page} shows a page in the web app from the moment it is
listed until it is deleted, including undo, failures and reloads.
Figure~\ref{fig:state-session} shows a sign-in session: it ends only when
Cognito refuses the refresh token, so a refresh that is cancelled by a reload,
or cannot reach Cognito, keeps the user signed in.

\begin{figure}[H]
\centering
\begin{tikzpicture}[>={Stealth[length=1.8mm]}]
  \node[startdot] (s) at (-2.1,0) {};
  \node[ustate] (shown)    at (0,0)     {Shown};
  \node[ustate] (pending)  at (5.4,0)   {Pending delete\\\tiny hidden, Undo toast};
  \node[ustate] (deleting) at (5.4,-2.9) {Deleting\\\tiny request sent};
  \node[ustate] (deleted)  at (10.6,-2.9) {Deleted};
  \node[enddot] (e) at (10.6,-4.6) {};

  \draw[trans] (s) -- node[tlabel, above] {listed} (shown);
  \draw[trans] (shown) -- node[tlabel, above=1mm] {delete / hide, show toast} (pending);
  \draw[trans] (pending.north) to[bend right=25] node[tlabel, above] {undo / show again} (shown.north);
  \draw[trans] (pending) -- node[tlabel, right] {6\,s elapsed, or tab closed / send DELETE} (deleting);
  \draw[trans] (deleting) -- node[tlabel, above] {204 or 404} (deleted);
  \draw[trans] (deleting.west) to[bend left=20] node[tlabel, sloped, below] {request failed / error toast} (shown.south);
  \draw[trans] (deleting.south west) .. controls ++(-0.4,-1.1) and ++(0.4,-1.1) .. node[tlabel, below] {reload [flushed $<$ 30\,s] / stay hidden, re-send} (deleting.south east);
  \draw[trans] (deleted) -- node[tlabel, right] {tab closed} (e);
  \draw[trans] (deleted.north) |- node[tlabel, pos=0.75, above] {same URL saved again [list fetched after delete]} ($(shown.north)+(0.4,1.55)$) -- ($(shown.north)+(0.4,0)$);
\end{tikzpicture}

\caption{State chart: a page in the web app (delete with undo)}
\label{fig:state-page}
\end{figure}

\begin{figure}[H]
\centering
\begin{tikzpicture}[>={Stealth[length=1.8mm]}]
  \node[startdot] (s) at (-2.0,0) {};
  \node[ustate] (out) at (0, 0)    {Signed out};
  \node[ustate] (sin) at (5.2, 2.3) {Signing in\\\tiny at Cognito};
  \node[ustate] (in)  at (10.4, 0)  {Signed in\\\tiny access token valid};
  \node[ustate] (ref) at (5.2,-2.3) {Refreshing\\\tiny token request};

  \draw[trans] (s) -- node[tlabel, above=1mm, pos=0.55] {[no session]} (out);
  \draw[trans] (s) |- node[tlabel, pos=0.8, above] {[refresh token stored]} (10.4,3.35) -- (in.north);
  \draw[trans] (out.north east) -- node[tlabel, sloped, above] {sign in / PKCE redirect} (sin.west);
  \draw[trans] (sin.south west) to[bend left=25] node[tlabel, sloped, below] {error or state mismatch} (out.east);
  \draw[trans] (sin.east) -- node[tlabel, sloped, above] {code exchanged / store tokens} (in.north west);
  \draw[trans] (in.west) -- node[tlabel, above] {sign out / clear session} (out.east |- in.west);
  \draw[trans] (in.south) |- node[tlabel, pos=0.75, below] {token near expiry, or reload} (ref.east);
  \draw[trans] (ref.north east) -- node[tlabel, sloped, above, pos=0.42] {new tokens, or unreachable\\/ keep session} (in.south west);
  \draw[trans] (ref.west) -| node[tlabel, pos=0.25, below] {400 or 401 / clear session} (out.south);
\end{tikzpicture}

\caption{State chart: a sign-in session in the web app}
\label{fig:state-session}
\end{figure}
